%
% Funktionenkörper
%
\subsection{Funktionenkörper}
In Kapitel~\ref{chapter:ratint} wurden das Problem studiert,
rationale Funktionen zu integrieren.
Die meisten Operationen bei der Lösung dieses Problems ergaben
wieder rationale Funktionen.
Nur ganz am Schluss traten Terme auf, deren Integrale mit dem
Logarithmus eines Polynoms berechnet werden können.
Dies zeigt, dass wir zum Aufbau von Stammfunktion zumindest die
arithmetischen Grundoperationen auf beliebige Funktionen anwenden können
müssen.
Die natürliche Struktur, in der sich die Integrationstheorie
abspielt, ist daher ein Körper von Funktionen.
Im Kapitel~\ref{chapter:ratint} war dies der Körper der rationalen
Funktionen.
In diesem Abschnitt muss es also darum gehen, diesen Funktionenkörper
massvoll zu erweitern, dass möglichst viele praktisch wichtige
Funktionen darin vorkommen.

%
% Koeffizientenkörper
%
\subsubsection{Koeffizientenkörper}
In Abschnitt~\ref{buch:ratint:section:bernoulli} haben wir als
Funktionenkörper den Körper $\mathbb{C}(x)$ verwendet, dann
aber später eingesehen, dass der Körper $\mathbb{C}$ zu gross ist.
Wenn man mit irrationalen Zahlen algebraisch rechnen muss, muss
man deren algebraischen Eigenschaften kennen.
Die Nullstellen eines Polynoms, die für die Partialbruchzerlegung
benötigt wurden, sind algebraische Element über einem Körper, der
alle Koeffizienten des Nennerpolynoms enthält.
Der Ausgangskörper für die Konstruktion der rationalen Funktionen
entsteht also durch Hinzufügen weiterer Konstanten zu den rationalen
Zahlen.
Die zusätzlichen Konstanten können algebraisch sein, wie die
Wurzel $\!\sqrt{2}$, oder transzendent, wie $\pi$, für die es keine
polynomielle Relation gibt.
Als Zahlenkörper, aus dem die Koeffizienten der Polynome stammen
muss daher immer eine Erweiterung von $\mathbb{Q}$ verwendet werden,
die durch Hinzufügen endlich vieler Elemente $c_1,\dots,c_n$ entstehen.
In der Kette
\[
\underbrace{\mathbb{Q}}_{\displaystyle = K_0}
\subset
\underbrace{\mathbb{Q}(c_1)}_{\displaystyle = K_1}
\subset
\underbrace{\mathbb{Q}(c_1,c_2)}_{\displaystyle = K_2}
\subset
\underbrace{\mathbb{Q}(c_1,c_2,c_3)}_{\displaystyle = K_3}
\subset
\ldots
%\subset
%\underbrace{\mathbb{Q}(c_1,c_2,c_3,\ldots,c_{n-1})}_{\displaystyle = K_{n-1}}
\subset
\underbrace{\mathbb{Q}(c_1,c_2,c_3,\ldots,c_{n-1},c_n)}_{\displaystyle = K_n}
=
K
\]
wird in jedem schritt genau ein neues Element
$c_k$ hinzugefügt, welches transzendent oder algebraisch ist über
dem vorangegangenen Körper.
Falls $c_k$ algebraisch ist, dann gibt es ein Minimalpolynom
$m_k(x) \in K_{k-1}[x]$, für welches $m_k(c_k)=0$ gilt.

Zu den transzendenten Elementen in $K$ gehören insbesondere auch
Parameter in den Ausdrücken.
Wenn ein CAS den quadratischen Ausdruck $ax^2+bx+c$ integrieren soll,
ohne den Koeffizienten $a$, $b$ und $c$ feste Werte zuzuordnen, dann
müssen diese Elemente als transzendent über jedem zugrundeliegenden
Körper betrachtet werden.
Damit sagt man, dass keine weiteren Beziehungen zwischen diesen
Koeffizienten bekannt sind.

%
% Das Funktionsargument
%
\subsubsection{Das Funktionsargument}
In einem Funktionsausdruck spielt ein Symbol die spezielle Rolle
des Funktionsarguments.
Aus algebraischer Sicht verhält sich der Platzhalter für das
Funktionsargument nicht anders wie jeder andere Parameter.
Im Ausdruck $f(x) = ax^2 + bx + c$ werden die Zeichen $a$, $b$ und $c$
wie gewöhnliche Zahlen behandelt, während das Symbol $x$ für das
Funktionsargument steht, für das man nach der Intuition, die man
beim ersten Kontakt mit dem Funktionsbegriff kennenlernt, beliebige
Werte einsetzen kann.
Aus algebraischer Sicht ist $x$ einfach nur ein weiteres transzendentes
Element, um das der Koeffizientenkörper erweitert wird.
Eine besondere Bedeutung erhält $x$ erst, wenn wir, wie das im
nächsten Abschnitt geschehen soll, mit der Ableitung die Konstanten
von der Variablen $x$ unterscheiden können.

\begin{beispiel}
Die rationale Funktion
\[
f(x) 
=
\frac{ax^2 + \!\sqrt{2}x - \sqrt{\pi}}{x^2+\pi}
\]
hat im Zählerpolynom die drei Koeffizienten $a$, $\!\sqrt{2}$ und
$\!\sqrt{\pi}$, im Nenner kommt ausserdem der Koeffizient $\pi$ hinzu.
Die Elemente $a$ und $\pi$ sind transzendent über den rationalen
Zahlen.
Die Elemente
$\!\sqrt{2}$
und
$\!\sqrt{\pi}$
entstehen
durch algebraische Erweiterung von $\mathbb{Q}$ mit dem
Minimalpolynom $m(x)=x^2-2$ bzw.~von $\mathbb{Q}(\pi)$ mit
dem Minimalpolynom $m(x)=x^2-\pi$.
Die Funktion ist daher ein Element im Körper der rationalen
Funktionen
\[
f(x)
\in
\mathbb{Q}(\!\sqrt{2},\pi,\!\sqrt{\pi},a,x).
\]
Man könnte auch versuchen, $\!\sqrt{\pi}$ als transzendentes Element
über $\mathbb{Q}$ zu betrachten.
Dann müsste man aber das Element $\pi$ als das Quadrat von $\!\sqrt{\pi}$
betrachten.
In diesem Fall wäre es zweckmässiger, die Funktion als
\[
f(x)
=
\frac{ax^2+\!\sqrt{2}-c}{x^2+c^2}
\]
mit einem transzendenten Element $c$ zu schreiben, welches für
$\sqrt{\pi}$ steht.
Diese Version macht es aber schwieriger, eine Beziehung zwischen $\pi$ 
und den trigonometrischen Funktionen herzustellen.
\end{beispiel}

Der kleinste Funktionenkörper, der die Konstanten in $K$ und das
Funktionsargument $x$ enthält, ist der Körper $K(x)$ der rationalen
Funktionen mit Koeffizienten in $K$.
Wir bezeichnen ihn auch mit $F=K(x)$.
Im Folgenden verwenden wir die Konvention, Funktionenkörper mit $F$
und Konstantenkörper mit $K$ zu bezeichnen.

